%% kapa: exam=1 round=21 year=2010 period=2 form=A subject=appraisal_law subject_ko=부동산관계법규 seq=13 number=53
%% source: Q-Net 감정평가사 기출문제 (cst00309) artlSeq=5024181 file=extracted/1008/5024181/2010년 제21회 감정평가사 1차시험 기출문제/2교시.hwp
%% conversion: kordoc 4.24.51 -> tools/gen1.py (auto); review=proofread-gpt-sol; flags=none
\begin{question}{53}
국토의 계획 및 이용에 관한 법령상 용적률의 최대한도가 같은 것끼리 묶어 놓은 것은? (단, 조례는 고려하지 않음)
\begin{choices}
  \item 제2종전용주거지역 - 계획관리지역 - 자연녹지지역
  \item 보전녹지지역 - 보전관리지역 - 계획관리지역
  \item 준주거지역 - 준공업지역 - 근린상업지역
  \item 제1종전용주거지역 - 자연녹지지역 - 생산녹지지역
  \item 전용공업지역 - 제3종일반주거지역 - 유통상업지역
\end{choices}
\end{question}
\answer{4}
